\chapter{THỰC NGHIỆM}

\section{Giới thiệu Dataset}

\subsection{Ali-CCP (Alibaba Click-Conversion Prediction)}
% TODO: Mô tả chi tiết dataset.
% Gợi ý:
%   - Nguồn: Alibaba Tianchi Platform.
%   - Quy mô: Số mẫu, số người dùng, số sản phẩm, thời gian thu thập.
%   - Cấu trúc: user_id, item_id, timestamp, label (click/purchase), device features.
%   - Tại sao chọn: Dataset thực tế từ chính Alibaba - đơn vị tiên phong DCCL.
%   - Thống kê cơ bản (bảng hoặc biểu đồ): Phân phối rating, sparsity, ...

\subsection{Tiền xử lý dữ liệu}
% TODO: Mô tả các bước tiền xử lý.
% Gợi ý:
%   - Lọc người dùng/sản phẩm có ít tương tác (K-core filtering).
%   - Chia tập train/validation/test (theo thời gian - temporal split).
%   - Xây dựng đồ thị tương tác: Node = User + Item, Edge = Interaction.
%   - So sánh dung lượng: Dữ liệu thô vs. Embeddings (minh họa lợi ích DCCL).

\section{Sơ đồ thiết kế thực nghiệm}

\subsection{Kiến trúc hệ thống tổng thể}
% TODO: Chèn hình kiến trúc hệ thống DCCL.
% Gợi ý nội dung hình:
%   - Cloud Server: GNN Training (GraphSAGE) → Item Embedding Store → Candidate API.
%   - Device Client: Local Subgraph → TFLite Ranker → Recommendations.
%   - Kênh giao tiếp: REST API (JSON).

% \begin{figure}[h!]
%     \centering
%     \includegraphics[width=0.85\textwidth]{figures/system_architecture.png}
%     \caption{Kiến trúc tổng thể hệ thống Device-Cloud Collaborative Learning}
%     \label{fig:architecture}
% \end{figure}

\subsection{Lưu đồ giải thuật}
% TODO: Vẽ lưu đồ minh họa quy trình inference end-to-end.
% Gợi ý (dùng tikz hoặc chèn hình):
%   Bước 1: Device gửi {user_id, k} lên Cloud.
%   Bước 2: Cloud thực hiện 3-stage retrieval (Graph → Vector → Popularity).
%   Bước 3: Cloud trả về {candidate_ids, embeddings}.
%   Bước 4: Device ranking bằng TFLite model.
%   Bước 5: Trả kết quả Top-K cho người dùng.

\section{Các phương pháp đánh giá}

% TODO: Mô tả 3-4 metrics sẽ sử dụng.

\subsection{Precision@K}
% Công thức: P@K = (số items liên quan trong top-K) / K

\subsection{Recall@K}
% Công thức: R@K = (số items liên quan trong top-K) / (tổng số items liên quan)

\subsection{NDCG@K (Normalized Discounted Cumulative Gain)}
% Công thức: NDCG@K = DCG@K / IDCG@K
% Trong đó: DCG@K = \sum_{i=1}^{K} \frac{rel_i}{\log_2(i+1)}

\subsection{Hit Rate@K}
% Công thức: HR@K = 1 nếu item liên quan xuất hiện trong top-K, ngược lại 0.

\section{Mô tả thực nghiệm}

\subsection{Thiết lập thực nghiệm}
% TODO: Mô tả cấu hình thực nghiệm.
% Gợi ý:
%   - Phần cứng: CPU/GPU dùng để huấn luyện Cloud model, thiết bị mô phỏng Device.
%   - Thư viện: PyTorch Geometric (PyG), TensorFlow Lite, FastAPI.
%   - Siêu tham số: Số lớp GNN, kích thước embedding, learning rate, batch size, epochs.
%   - Baseline so sánh: BPR-MF, LightGCN, PureSVD.

\subsection{Kết quả thực nghiệm}
% TODO: Trình bày kết quả dưới dạng bảng so sánh.

\begin{table}[h!]
    \centering
    \caption{Kết quả so sánh hiệu năng giữa các phương pháp trên tập test}
    \renewcommand{\arraystretch}{1.4}
    \begin{tabularx}{\textwidth}{|l|X|X|X|X|}
        \hline
        \textbf{Phương pháp} & \textbf{Precision@10} & \textbf{Recall@10} & \textbf{NDCG@10} & \textbf{HR@10} \\ \hline
        BPR-MF       & --    & --    & --    & -- \\ \hline
        LightGCN     & --    & --    & --    & -- \\ \hline
        DCCL (Nhóm)  & --    & --    & --    & -- \\ \hline
    \end{tabularx}
    \label{tab:results}
\end{table}

\subsection{Phân tích kết quả}
% TODO: Nhận xét và phân tích kết quả từ bảng so sánh.
% Gợi ý:
%   - So sánh hệ thống DCCL của nhóm với các baseline.
%   - Phân tích ảnh hưởng của k (số ứng viên) đến chất lượng khuyến nghị.
%   - Vẽ biểu đồ: Bar chart hoặc line chart thể hiện Precision/NDCG theo K.

% \begin{figure}[h!]
%     \centering
%     \includegraphics[width=0.75\textwidth]{figures/ndcg_comparison.png}
%     \caption{So sánh NDCG@K của các phương pháp}
%     \label{fig:ndcg}
% \end{figure}

\subsection{Phân tích chi phí truyền thông}
% TODO: So sánh băng thông cần thiết giữa truyền dữ liệu thô và truyền embeddings.
% Gợi ý (vẽ biểu đồ cột):
%   - Truyền dữ liệu thô: X MB/request
%   - Truyền embeddings: Y KB/request
%   => Tiết kiệm Z lần băng thông.
